\section{Assembly of the dyadic estimate}\label{sec:assembly}

We prove the modulated delayed Haar estimate stated in
\Cref{prop:dyadic}. First we bound the contribution of a single plateau,
using its representatives to recover each actual modulation. We then sum
these contributions with the plateau starting weights, keeping the
live Fisher and edge charges global.

\begin{proof}[Proof of \Cref{prop:dyadic}]
The zero input is immediate.  For positive input, apply the finite
procedure of \Cref{sec:procedure}, with the common stopping rule in the
statement.  In particular, the cutoff always refers to the original
$\rho$, including when a subpiece is being processed.  The parameter
constants are fixed in the order specified there, and $d_0$ is increased
to satisfy their finitely many large-$d$ requirements.

\subsection*{The contribution of one plateau}

For a plateau $e$, write $p_e$ for its root, $\alpha_e$ for its
normalization, and $\mathbf U_e$ for uniform probability on $p_e$.
On an output path $x\in p_e$, let $\mathcal P_e(x)$ be its consecutive
jumping nodes belonging to this plateau.  The post density at a node
$p\in\mathcal P_e(x)$ is denoted by $a_p$.  Thus all removals at $p$
have already occurred, and $a_p$ is the density whose two prospective
child moments are used in that jump.  In sums over nodes $p$, write
$j=\operatorname{depth}(p)$.  Define
\[
 c_p^u=\frac1{\alpha_e}\E_p[a_p(Y)\chi_u(Y)r_{j+1}(Y)],
 \qquad
 \mathcal T_e^u(x)=\sum_{p\in\mathcal P_e(x)}
           c_p^u\epsilon_j(x)r_{j+2}(x).
\]
At each jumping node the post pieces partition the original density.
Consequently their transforms satisfy the exact identity
\[
 \mathcal H_{I,J}^{u}\rho(x)
    =\sum_{e:\,x\in p_e}\alpha_e\mathcal T_e^u(x).
\]
Put $w_e=\mathbf U(p_e)\alpha_e$. The triangle inequality and integration
with respect to $\mathbf U$ give
\begin{equation}\label{eq:09-plateau-reduction}
 \int_I\max_{\substack{u\in\mathbb Z\\|u|\le Q}}
       |\mathcal H_{I,J}^{u}\rho|\,d\mathbf U
 \le\sum_e w_e\mathbb E_{\mathbf U_e}
                     \sup_{u\in\mathbb Z}|\mathcal T_e^u|.
\end{equation}
We will bound this weighted sum by $Cd$, charging the live Fisher and
edge terms across all pieces and the remaining terms one plateau at a time.

The local mass bound gives $|c_p^u|\le2$ and prospective child masses
at most $4$, in this normalization.  Let $K_e(x)$ denote the sum of all
normalized removal masses $\xi_p$ on the plateau's processing path,
including the intervening and terminal removals allowed in
\Cref{prop:capture}; then
\begin{equation}\label{eq:09-kills}
 \E_{\mathbf U_e}K_e\le1.
\end{equation}
For live plateaus set $d_1=d$, and for static plateaus set
$d_1=d+2+\log^-\alpha_e$.

Use the main levels
$\mathcal B_{\rm m}=\{2^{-j}:b_0\le2^{-j}\le1\}$ on live plateaus.
Use scalar levels
$\mathcal B_{\rm s}(B)=\{B2^{-j}:j\ge0\}$, where $B=b_0$ on a live
plateau and $B=1$ on a static plateau.  At the shared level $b_0$ the
main and scalar constructions are distinct.  By
\Cref{lemma:paths-simultaneous}, they have nonnegative, frequency-independent
path parameters $t_b^{\rm m},t_b^{\rm s}$ and miss-block lengths
$H_b^{\rm m},H_b^{\rm s}$, satisfying
\begin{align}
 \E_{\mathbf U_e}t_b^{\rm m}
   &\le C k_{b/4}+C\log(Cd/b),
       &&b\in\mathcal B_{\rm m},\label{eq:09-main-parameters}\\
 \E_{\mathbf U_e}t_b^{\rm s}
   &\le C(d_1+\log(1/b)),
       &&b\in\mathcal B_{\rm s}(B),\label{eq:09-scalar-parameters}\\
 \E_{\mathbf U_e}H_b^{\rm m},\quad
 \E_{\mathbf U_e}H_b^{\rm s}
   &\le C\log(Cd_1/b)+C\log(1/\sigma)
       &&\text{for the levels where they are used.}
       \label{eq:09-miss-parameters}
\end{align}
These statements already include all snapshot births, representative
choices, starts, checkpoints, and dyadic list caps.  In particular we may
use the same parameters for every actual integer $u$.

At a live jumping piece write
\[
 E_{p,b}^{\rm m}=e_{p,b}\mathbf1_{\{p\text{ bad at main level }b\}}.
\]
For the scalar construction let $\bar b(p)$ be its largest bad level,
if there is one, and put
\[
 E_p^{\rm s}=e_{p,\bar b(p)}
 \quad\text{if a bad scalar level exists},\qquad E_p^{\rm s}=0
 \quad\text{otherwise}.
\]
A nonempty subset of the dyadic levels bounded above by $B$ has a
largest member.  This choice is determined by the jumping piece before
its bit and is independent of $u$.  The scalar edge regions are nested
as $b$ increases: their defining $l_b$ decreases.  Hence
\begin{equation}\label{eq:09-edge-monotonicity}
 b\le\bar b(p)\quad\Longrightarrow\quad e_{p,b}\le e_{p,\bar b(p)}.
\end{equation}

\subsection*{Main amplitudes and exact decoding}

Fix an integer $u$ and a live plateau path.  Starting with its first main
activation, including that jump, retain the largest main level activated
so far.  This level is nondecreasing.  Its positive intervals therefore
use each level $b$ at most once.  On the interval assigned to $b$, keep
the representative $s$ and decomposition supplied at its first
activation:
\begin{equation}\label{eq:09-main-decomposition}
 u=s+n+a',\qquad \Delta_A(n)\le r_{\rm sep},\qquad
 |a'|\,|p_{\rm act}|\le C\sigma.
\end{equation}
The form $A$ is fixed on this plateau.  All subsequent nodes in the
interval lie below $p_{\rm act}$.

For clarity, the comparison of criteria uses no loss depending on the
number of representatives.  For a prospective child $q$ of $p$, write
$f_q^v(z)=\alpha_e^{-1}\E_q[a_p(Y)\chi_v(Y)K_A(z-Y)]$ and
$M_q(z)=\alpha_e^{-1}\E_q[a_p(Y)K_A(z-Y)]$.
For the fixed output point $x\in p$, conditional Cauchy--Schwarz gives
\[
 \left\|f_q^u-\chi_n(\cdot)\chi_{a'}(x)f_q^s
       \right\|_{L^2(M_q^{-1}dz)}
 \le C\Delta_A(n)+C|a'|\,|p|.
\]
Indeed the square of the channel part is at most
\[
 \frac1{\alpha_e}\E_q\left[
     a_p(Y)\int |\chi_n(Y)-\chi_n(z)|^2K_A(z-Y)\,dz\right]
 =\frac{\E_q a_p}{\alpha_e}\Delta_A(n)^2,
\]
and the spatial part uses
$|\chi_{a'}(Y)-\chi_{a'}(x)|\le2\pi|a'|\,|p|$ and
$\E_q a_p/\alpha_e\le4$.  The triangle inequality in this weighted
Hilbert space proves
\begin{equation}\label{eq:09-criterion-comparison}
 \left|\sqrt{G_q^u}-\sqrt{G_q^s}\right|
 \le C(r_{\rm sep}+\sigma)\ll b_0\le b.
\end{equation}

If $b<1$, level $2b$ has not activated at any jump of this
interval.  Its non-bad high jumps before activation lie in one block of
at most $H_{2b}^{\rm m}$ depths, by \Cref{prop:capture}.  Remove the
intersection of that entire block with the current interval and pay
$C H_{2b}^{\rm m}$ directly, using $|c_p^u|\le2$.  The remaining jumps
form at most two contiguous intervals.  On each, a nonexceptional jump
is not high for $u$ at $2b$; thus
\Cref{eq:09-criterion-comparison} makes it inner-valid for $s,b$ in
\Cref{prop:vector-path}.  Here the exceptional jumps are precisely the
bad main jumps at $2b$.  When $b=1$, every jump is inner-valid after the
fixed constant in that definition is chosen large enough, since all
$G_q^s\le4$; no block is removed and no exception is needed.

We next verify the decoding of the vector estimate back to the actual
scalar modulation.  Put
\[
 \mathfrak a_p^s(z)
   =\frac1{\alpha_e}\E_p[a_p(Y)\chi_s(Y)r_{j+1}(Y)K_A(z-Y)],
 \qquad c_A(n)=\int\chi_n(\theta)K_A(\theta)\,d\theta.
\]
Evenness and positivity of the thermal Fourier coefficients imply
\begin{equation}\label{eq:09-decoding}
 c_A(n)=1-\tfrac12\Delta_A(n)^2\ge\tfrac12,
 \qquad
 \int\chi_n(z)\mathfrak a_p^s(z)\,dz=c_A(n)c_p^{s+n}.
\end{equation}
This is an exact identity, obtained by putting $z=Y+\theta$ in the
inner integral.  On any subset of the nodes below $p_{\rm act}$,
\[
 \sum_p|c_p^u-\chi_{a'}(x)c_p^{s+n}|
 \le C|a'|\sum_p|p|
 \le C|a'|\,|p_{\rm act}|\le C\sigma.
\]
The sum of lengths is geometric because the nodes lie on one descending
path.  The bound applies in particular to each retained subinterval.
Since the multipliers are common to all label coordinates,
\Cref{eq:09-decoding} now gives
\begin{equation}\label{eq:09-vector-to-scalar}
 \left|\sum_{p\in L}c_p^u\epsilon_jr_{j+2}\right|
 \le2\int\left|\sum_{p\in L}\mathfrak a_p^s(z)
                           \epsilon_jr_{j+2}\right|\,dz+C\sigma
\end{equation}
for every retained subinterval $L$.

Apply \Cref{prop:vector-path} to $L$.  Its Fisher sum is charged only
on $L$.  The intervals at distinct retained levels are disjoint, and
there are at most two retained subintervals at each level.  Thus these
Fisher charges have an absolute bounded multiplicity at every jumping
piece.  All kill and exceptional sums may use the whole plateau path.
Writing the exceptional proxy as
$e'_p=C b_0^{-2}E_{p,2b}^{\rm m}$, and using
$b^{-1/2}b_0^{-2}\le b_0^{-3}$ for $b\ge b_0$, we have bounded the whole
main part by a constant times
\begin{equation}\label{eq:09-main-path-bound}
 \sum_pJ_p+b_0^{-3}\sum_p\sum_{b\in\mathcal B_{\rm m}}E_{p,b}^{\rm m}
 +\sum_{b\in\mathcal B_{\rm m}}
       \left[b^{-1/2}(1+K_e)+b^{1/2}t_b^{\rm m}+H_b^{\rm m}\right].
\end{equation}
The constants in the vector estimates and the geometric decoding errors
are absorbed by the displayed $b^{-1/2}$ terms.

\subsection*{The scalar prefix and all static paths}

On a live plateau consider the prefix strictly before the first main
activation, or the whole path if no main activation occurs.  On a static
plateau consider the entire path.  Use the scalar levels through $B$
specified above, again retaining the largest activated level so far,
including activations at the current jump, and using zero if none has
activated.  There are finitely many changes
on this finite path even though the available scalar levels extend
arbitrarily close to zero.  Each positive level $b$ is retained on at
most one contiguous interval.

For its fixed representative $s$, scalar activation supplies
$u=s+a'$ with $|a'|\,|p_{\rm act}|\le Cb^4$, with the sufficiently
small constant in \Cref{prop:capture}.  On this level interval sum first
over the fixed scalar mask
\[
 \max\{|f_{p,+}^s|,|f_{p,-}^s|\}\le C_3b.
\]
This is precisely the masked transform in \Cref{prop:scalar-path}; that
proposition applies to the masked sum over the whole contiguous level
interval.  Freezing $\chi_{a'}$ at $x$ changes this masked sum by at most
$Cb^4$, by the same geometric sum of interval lengths as above.  Hence
all such masked sums cost at most
\begin{equation}\label{eq:09-scalar-mask-sum}
 C\sum_{b\in\mathcal B_{\rm s}(B)}b^{1/4}
                         (1+K_e+t_b^{\rm s}).
\end{equation}

Treat every remaining jump directly.  This includes all jumps while the
retained level is zero.  Put
$A'_p=\max\{|f_{p,+}^u|,|f_{p,-}^u|\}\le4$ in the unlabelled scalar
notation.  The direct jump costs at most $A'_p$, and a zero value needs
no charge.  If a positive level $b$ is currently retained, the jump is
outside its mask; the spatial phase comparison gives
$A'_p\ge C_3b-Cb^4>4b$, on taking $C_3$ large enough.

For a live prefix, any direct jump with $A'_p>4b_0$ is high at main
level $b_0$.  In fact $f_{p,\pm}^u=\int f_{p,\pm}^u(z)\,dz$ and
Cauchy--Schwarz gives
$|f_{p,\pm}^u|^2\le(\int M_{p,\pm})G_{p,\pm}^u\le4G_{p,\pm}^u$.
No main level has yet activated in this prefix.  All its non-bad main
high jumps at $b_0$ are therefore paid by its one miss block, and its
bad jumps are paid by
$C b_0^{-2}E_{p,b_0}^{\rm m}$, since each direct coefficient is at most
$4$.  These costs are already included, up to constants, in
\Cref{eq:09-main-path-bound}.

At every other nonzero direct jump we have $A'_p\le4B$; this holds
automatically on a static path with $B=1$.  Choose the largest available
scalar level $\widetilde b\le A'_p$.  Then
\[
 A'_p/4\le\widetilde b\le A'_p,
\]
and the jump is high at $\widetilde b$.  This level has not activated:
if the current level is positive, $A'_p>4b$ implies
$\widetilde b>b$, and otherwise none has activated.  All non-bad jumps
assigned to a fixed $\widetilde b$ consequently lie in its single
preactivation miss block and cost at most
$C\widetilde b H_{\widetilde b}^{\rm s}$ together.  If the assigned
level is bad, then
$e_{p,\widetilde b}\ge c\widetilde b$ and
\Cref{eq:09-edge-monotonicity} gives
\[
 A'_p\le4\widetilde b\le C e_{p,\widetilde b}\le C E_p^{\rm s}.
\]
Thus only one predictably selected scalar edge level per jump is
charged.  This establishes the scalar accounting for every positive
amplitude, including arbitrarily small amplitudes.

Combining these conclusions, the same frequency-independent quantity
bounds $\sup_u|\mathcal T_e^u|$ on a live plateau:
\begin{align}
 C^{-1}\sup_{u\in\Z}|\mathcal T_e^u|
 \le{}&\sum_{p\in\mathcal P_e(x)}
       \left(J_p+b_0^{-3}\sum_{b\in\mathcal B_{\rm m}}E_{p,b}^{\rm m}
                     +E_p^{\rm s}\right)\notag\\
 &+\sum_{b\in\mathcal B_{\rm m}}
       \left[b^{-1/2}(1+K_e)+b^{1/2}t_b^{\rm m}+H_b^{\rm m}\right]\notag\\
 &+\sum_{b\in\mathcal B_{\rm s}(b_0)}
       \left[b^{1/4}(1+K_e+t_b^{\rm s})+bH_b^{\rm s}\right].
       \label{eq:09-live-majorant}
\end{align}
For a static plateau the corresponding bound is
\begin{equation}\label{eq:09-static-majorant}
 \sup_{u\in\Z}|\mathcal T_e^u|
 \le C\sum_pE_p^{\rm s}
    +C\sum_{b\in\mathcal B_{\rm s}(1)}
       \left[b^{1/4}(1+K_e+t_b^{\rm s})+bH_b^{\rm s}\right].
\end{equation}
The path constructions and bounds used here hold simultaneously outside
one $\mathbf U_e$-null set.  The countable scalar-level sums have finite
expectation by the estimates below.  In particular no exceptional set
depending on the actual modulation is introduced in these majorants.

\subsection*{Summing levels and live weights}

Write $b_0=2^{-M}$, so $M=O(1+\log h)$.  Since
$k_b=A_0(d/h)(1+\log(1/b))^2$, the main cardinality tails satisfy
\begin{equation}\label{eq:09-main-series}
 \sum_{b\in\mathcal B_{\rm m}}b^{1/2}k_{b/4}
 =A_0\frac dh\sum_{j=0}^{M}2^{-j/2}
                          \bigl(1+(j+2)\log2\bigr)^2
 \le C\frac dh.
\end{equation}
The additional logarithms in \Cref{eq:09-main-parameters} sum to
$O(\log(Cd))$ with the same $b^{1/2}$ weights.  Also,
\[
 \sum_{b\in\mathcal B_{\rm m}}b^{-1/2}\E(1+K_e)
       \le Cb_0^{-1/2}\le Ch^8,
 \qquad
 \sum_{b\in\mathcal B_{\rm m}}\E H_b^{\rm m}
       \le C(1+\log h)\bigl(\log(Cd)+\log h\bigr).
\]
Every fixed polynomial in $h$, multiplied by $1+\log d$, is
$o(d/h)$ because $h=\log\log d$.  Thus the second line of
\Cref{eq:09-live-majorant} has expectation $O(d/h)$.

For the scalar series write $b=B2^{-j}$.  The convergent sums
$\sum_{j\ge0}2^{-j/4}$ and $\sum_{j\ge0}j2^{-j/4}$, together with
\Cref{eq:09-kills,eq:09-scalar-parameters,eq:09-miss-parameters}, give
\begin{align}
 &\sum_{b\in\mathcal B_{\rm s}(B)}
  \E\left[b^{1/4}(1+K_e+t_b^{\rm s})+bH_b^{\rm s}\right]\notag\\
 &\hspace{1cm}\le
 C B^{1/4}\bigl(d_1+1+\log(1/B)+\log(1/\sigma)\bigr).
 \label{eq:09-scalar-series}
\end{align}
For example the miss-block part alone is bounded by
$CB(\log(Cd_1/B)+\log(1/\sigma)+1)$, which is no larger than the
displayed right side.  On live plateaus $B=b_0\le h^{-16}$ and $d_1=d$,
so this is $O(d/h^4)$ and hence $O(d/h)$.  On static plateaus $B=1$ and
$d_1\ge d$, so $\log(1/\sigma)=40\log h\le Cd_1$ makes it $O(d_1)$.
The local scalar edge bound in \Cref{prop:capture} therefore proves
\begin{equation}\label{eq:09-static-plateau}
 \E_{\mathbf U_e}\sup_u|\mathcal T_e^u|\le C(d+2+\log^-\alpha_e)
 \quad\text{for every static plateau.}
\end{equation}

For every nonnegative jump quantity
$V_p$ of a plateau, change of the uniform root measure gives exactly
\begin{equation}\label{eq:09-weight-conversion}
 \sum_e w_e\E_{\mathbf U_e}\sum_{p\in\mathcal P_e(x)}V_p
   =\sum_{\text{jumping pieces }p}\mathbf U(p)\alpha_e V_p.
\end{equation}
The Fisher bound in \Cref{prop:live-budgets} bounds the right side by
$Cd$ when $V_p=J_p$ and $e$ is live.  The global scalar edge assertion
of \Cref{prop:capture} does the same for $V_p=E_p^{\rm s}$.  Its global
main edge assertion yields
\[
 b_0^{-3}\sum_{\text{live jumping pieces }p}
       \mathbf U(p)\alpha_e\sum_{b\in\mathcal B_{\rm m}}E_{p,b}^{\rm m}
 \le \poly(h)(h+\log d)=o(d).
\]
Finally $\sum_{e\text{ live}}w_e\le Ch$ by
\Cref{prop:live-budgets}.  Multiplying the $O(d/h)$ per-plateau costs
in \Cref{eq:09-live-majorant} by these weights therefore proves
\begin{equation}\label{eq:09-live-total}
 \sum_{e\text{ live}}w_e\E_{\mathbf U_e}\sup_u|\mathcal T_e^u|\le Cd.
\end{equation}

\subsection*{Static roots, small mass, and rare transfers}

Consider static handling starting with a positive piece of mean $m$ at
an interval $q$.  Its later plateaus restart only when their incoming
mean exceeds twice the previous normalization.  Along any input path
each new normalization is therefore more than twice its predecessor,
and every normalization is at least $m$.  A plateau with a jump has
normalization at most $e^{4h}$; allowing a final non-jumping overshoot
changes this bound only by a factor of two.  Consequently the number
$N_{\rm st}$ of static starts on the path obeys
\begin{equation}\label{eq:09-static-count}
 N_{\rm st}\le C(1+h+\log^-m),\qquad
 d+2+\log^-\alpha_e\le d+2+\log^-m.
\end{equation}
Starts with no subsequent jump may also simply be omitted.

Let $\mathbf P_q$ be the probability with density equal to this entry
piece divided by $m$ relative to uniform probability on $q$.  Since the
piece is never split again, the probability of reaching a later static
plateau root $p_e$ is
$\mathbf U_q(p_e)\alpha_e/m$.  Thus the sum of its starting weights is
the original weight $\mathbf U(q)m$ multiplied by
$\E_{\mathbf P_q}N_{\rm st}$.  Applying
\Cref{eq:09-static-plateau,eq:09-static-count}, the entire static
descendant family costs at most
\begin{equation}\label{eq:09-static-family}
 C\mathbf U(q)m(1+h+\log^-m)(d+2+\log^-m).
\end{equation}
This uses the input path only to sum weights; each individual transform
estimate remains an integral over its uniform output path.

If $0<m_{\rm top}<1/h$, the procedure takes the top static immediately.
Writing $\ell=\log(1/m_{\rm top})$, its total cost is bounded by
\[
 Cm_{\rm top}(1+h+\ell)(d+2+\ell)\le Cd.
\]
For an explicit verification, $m_{\rm top}h\le1$,
$\sup_{0<t\le1}t\log(1/t)=1/e$, and
$\sup_{0<t\le1}t\log^2(1/t)=4/e^2$.
Expanding the product and using $h\le d$ proves the inequality.

It remains to sum static families transferred out of an initially live
top, so $m_{\rm top}\ge1/h$.  Every true input path undergoes at most
one such transfer.  Let $A_{\rm tr}$ be the transfer event under its
input probability $\mathbf P$, and let $X_*$ be the maximum negative
logarithm of the mean through transfer or stopping, including the
transfer state.  \Cref{prop:live-budgets} gives
\begin{equation}\label{eq:09-transfer-data}
 \mathbf P(A_{\rm tr})\le Ch^{-10},\qquad
 \|X_*\|_{L^2(\mathbf P)}\le Ch.
\end{equation}
This includes both charge/count transfers and the small-mass flag.
In particular the small-mass event has $X_*>d$, hence probability at
most $Ch^2/d^2\le Ch^{-10}$.  The cap accounting and the evidence
argument in that proposition bound the probability of the other transfers;
there is no further sequence of live restarts after transfer.

For a transfer piece of mean $m$, its true reaching probability is
$\mathbf U(q)m/m_{\rm top}$ and $\log^-m\le X_*$ on that event.
Summing \Cref{eq:09-static-family} over all transfer roots therefore
bounds their entire cost by
\[
 Cm_{\rm top}\E_{\mathbf P}
       [\mathbf1_{A_{\rm tr}}(1+h+X_*)(d+2+X_*)].
\]
Only the second moment in \Cref{eq:09-transfer-data} is needed.  Indeed,
with $\delta=\mathbf P(A_{\rm tr})$,
\begin{align*}
 &\E[\mathbf1_{A_{\rm tr}}(1+h+X_*)(d+2+X_*)]\\
 &\quad=(1+h)(d+2)\delta
       +(d+3+h)\E[\mathbf1_{A_{\rm tr}}X_*]
       +\E[\mathbf1_{A_{\rm tr}}X_*^2]\\
 &\quad\le Cdh^{-9}+C(d+h)h^{-4}+Ch^2\le Cd.
\end{align*}
Here Cauchy--Schwarz gives
$\E[\mathbf1_{A_{\rm tr}}X_*]\le\sqrt\delta\,\|X_*\|_2\le Ch^{-4}$,
whereas the quadratic term uses its full bound
$\E X_*^2\le Ch^2$.  Multiplying by $m_{\rm top}\le2$ preserves $Cd$.

\subsection*{Completion}

The live estimate \Cref{eq:09-live-total} and the static-family estimates
bound the right side of \Cref{eq:09-plateau-reduction} by $Cd$. This proves
\Cref{eq:dyadic-estimate}.
Every construction was in fixed finite depth and every estimate was
independent of $Q$.  Increasing $Q$ and applying monotone convergence
proves the assertion for all integer modulations as well.
\end{proof}
